\section*{الفصل \ref{ch:b3:endocrinology} --- علم الغدد الصماء والاستتباب}

\begin{solution}{exo:b3:endocrinology:1}
\omterm{def:b3:endocrinology:glucose}{الأنسولين}: ببتيد، ومستقبل سطحي (كيناز تيروزين)، يعمل في
دقائق، وعمر نصفه نحو \qty{5}{min}. والكورتيزول: ستيرويد،
\omterm{def:b3:endocrinology:hormone}{ومستقبل نووي}، وساعات، وعمر نصفه \qty{80}{min} (وجلّه مرتبط
بحامل). والأدرينالين: أمين، ومستقبلات سطحية مقترنة بالبروتين
G، وثوانٍ، وعمر نصفه نحو \qty{2}{min}. والثيروكسين: حمض أميني
مُيَوْدن، \omterm{def:b3:endocrinology:hormone}{ومستقبل نووي}، وأيام، وعمر نصفه \qty{7}{d} (مرتبط
بحوامل). والفازوبريسين: ببتيد، ومستقبل سطحي مقترن بالبروتين
G (عبر cAMP في \omterm{def:b3:renal-osmoregulation:nephron}{القناة الجامعة})، ودقائق، وعمر نصفه نحو
\qty{15}{min}.
\end{solution}

\begin{solution}{exo:b3:endocrinology:2}
\omterm{def:b3:endocrinology:axes}{الوطاء} (TRH) $\to$ \omterm{def:b3:endocrinology:axes}{النخامى} (TSH) $\to$ الدرق (T$_{4}$)، مع
كبت T$_{4}$ للطبقتين العلويتين. والدرق المهدَّم: T$_{4}$
منخفض، وTSH مرتفع. \omterm{def:b3:cancer-biology:hallmarks}{والورم} المفرز \omterm{def:b3:endocrinology:hormone}{للهرمون} TSH: TSH مرتفع وT$_{4}$
مرتفع (إذ تخفق التغذية الراجعة في كبت \omterm{def:b3:cancer-biology:hallmarks}{الورم}). وأقراص
الثيروكسين الزائدة: T$_{4}$ مرتفع، وTSH مكبوت. وعوز اليود:
T$_{4}$ منخفض أو في أسفل السويّ، وTSH مرتفع، والغدة تنمو تحته
إلى دُراق.
\end{solution}

\begin{solution}{exo:b3:endocrinology:3}
الدماغ يموت من غلوكوز منخفض في دقائق ولا يؤذيه المرتفع إلا
على مدى سنين؛ والتطور لقي المجاعة أكثر بكثير مما لقي الولائم.
فالدفاع ضد الهبوط مكرَّر (\omterm{def:b3:endocrinology:glucose}{الغلوكاغون}، والأدرينالين،
والكورتيزول، \omterm{def:b3:endocrinology:hormone}{وهرمون} النمو) والدفاع ضد الارتفاع مفرد. والتنبؤ:
أن أنسولينًا أكثر مما ينبغي قاتل حادًّا (سُبات نقص سكر
الدم)، وأن أقل مما ينبغي داء مزمن --- وهو ما تراه العيادة.
\end{solution}

\begin{solution}{exo:b3:endocrinology:4}
\omterm{def:b3:endocrinology:clock}{المزامن} إشارة بيئية تُدرّج ساعةً: الضوء، وتوقيت الطعام،
والحرارة، والرياضة، والجدول الاجتماعي. والساعة الرئيسة
تستعمل الضوء، عبر الخلايا العقدية الحاملة للميلانوبسين في
\omterm{def:b3:sensory-systems:retina}{الشبكية}. وبعد ثماني مناطق زمنية لا تنزاح الساعة إلا نحو
ساعة في اليوم، فيسير النوم والكورتيزول والحرارة والشهية
أسبوعًا على التوقيت القديم بينما تنجرف الساعات المحيطية،
التي تعيد الوجبات ضبطها، على وتيرتها هي: أي وعكة الطيران،
وهي أسوأ شرقًا، إذ تقديم الساعة أصعب من تأخيرها.
\end{solution}

\begin{solution}{exo:b3:endocrinology:5}
نصف الاستجابة القصوى عند إشغال $250$ مستقبلًا: أي $H = K_{d}\times
250/(\num{20000} - 250) = \qty{1}{nM}\times 0.0127 = \qty{12.7}{pM}$،
أي أدنى من $K_{d}$ بثمانين ضعفًا. ومع \num{2000} مستقبل: $250/1750
\times \qty{1}{nM} = \qty{143}{pM}$، أي أقل حساسية بأحد عشر ضعفًا.
\omterm{def:b3:endocrinology:glucose}{فالأنسولين} المرتفع المطوَّل ينظّم مستقبلاته تنظيمًا نازلًا،
فيحتاج الأثر نفسه إلى \omterm{def:b3:endocrinology:glucose}{أنسولين} أكثر: وهو مكوّن من مكوّنات
مقاومة \omterm{def:b3:endocrinology:glucose}{الأنسولين}، وهي تتغذى على نفسها.
\end{solution}

\begin{solution}{exo:b3:endocrinology:6}
$T_{4} = 13$: $1.5\,\mathrm{e}^{0.92} = \qty{3.8}{mU/L}$؛ وعند $11$:
$1.5\,\mathrm{e}^{1.84} = \qty{9.4}{mU/L}$؛ وعند $9$: $1.5\,\mathrm{e}^{2.76}
= \qty{24}{mU/L}$؛ وعند $25$: $1.5\,\mathrm{e}^{-4.6} = \qty{0.015}{mU/L}$. وعند
\qty{11}{pmol/L} يكون T$_{4}$ لا يزال ضمن \qtyrange{10}{20}{} بينما
TSH أكثر من ضعف حده الأعلى (قصور درق دون سريري)؛ وعند $13$
يكون TSH داخل المجال بالكاد؛ وعند $9$ كلاهما شاذ؛ وعند $25$
كلاهما شاذ في الاتجاه الآخر.
\end{solution}

\begin{solution}{exo:b3:endocrinology:7}
$L = s\beta/(a\gamma) = 0.2\times 0.05/(0.02\times 0.1) = 5$. والفائض
المستقر $g^{*} = (u/a)/(1 + L) = (2/0.02)/6 = \qty{16.7}{mg/dL}$؛ وبلا
تغذية راجعة $u/a = \qty{100}{mg/dL}$. وفائض \omterm{def:b3:endocrinology:glucose}{الأنسولين} $i^{*} = \beta g^{*}/
\gamma = 0.05\times 16.7/0.1 = \qty{8.3}{mU/L}$.
\end{solution}

\begin{solution}{exo:b3:endocrinology:8}
المميز $(a - \gamma)^{2} - 4s\beta = 0.0064 - 0.04 < 0$:
فالرجوع تذبذبي. و$\omega = \sqrt{a\gamma + s\beta - (a+\gamma)^{2}/4} =
\sqrt{0.002 + 0.01 - 0.0036} = \qty{0.092}{min^{-1}}$، والدور $2\pi/
\omega = \qty{68}{min}$؛ وتهبط السعة بعامل $\mathrm{e}$ في $2/(a +
\gamma) = \qty{16.7}{min}$. وبتنصيف $s$: $s\beta = 0.005$، والمميز
$0.0064 - 0.02 < 0$، فلا يزال تذبذبيًا، و$\omega = \sqrt{0.007 - 0.0036}
= \qty{0.058}{min^{-1}}$، والدور \qty{108}{min}؛ وزمن الخمود بلا
تغيّر؛ و$L = 2.5$. فعروة أضعف ترجع أبطأ وتحمل خطأً مستقرًا أكبر.
\end{solution}

\begin{solution}{exo:b3:endocrinology:9}
ظل البريدنيزولون سنةً يكبت CRH وACTH؛ فضمر قشر الكظر إذ لم
يُحرَّض. وعند التوقف لا ستيرويد خارجيًا، ويحتاج \omterm{def:b3:endocrinology:axes}{الوطاء}
\omterm{def:b3:endocrinology:axes}{والنخامى} إلى أسابيع ليستأنفا، والكظر الضامر ما كان ليستجيب
\omterm{def:b3:endocrinology:hormone}{لهرمون} ACTH لو جاء. والخمج يطلب عشرة أضعاف الكورتيزول السويّ؛
فإذا لم يكن منه شيء، توسعت الأوعية، وهبط الضغط والغلوكوز:
أي سَورة كظرية. والمقيس: ACTH منخفض (فالآفة مركزية)،
والكورتيزول منخفض، وارتفاع كورتيزولي رديء لحقن ACTH (إذ ضمرت
الغدة) --- بخلاف داء أديسون، حيث يكون ACTH مرتفعًا. ومن ثم
كان الخفض التدريجي البطيء.
\end{solution}

\begin{solution}{exo:b3:endocrinology:10}
بمتغير واحد: $\dot x = f(x)$ عند $f' < 0$ له نقطة ثابتة واحدة
$x^{*}$؛ فعند $x > x^{*}$ يكون $\dot x < 0$ وينقص $x$ نقصًا
رتيبًا نحو $x^{*}$، ولا يعبره أبدًا (لوحدانية الحلول)،
وبالتناظر تحته: فلا تذبذب. وبمتغيرين: تقدر المصفوفة على أن
تكون لها قيم ذاتية عقدية، كما في المبرهنة، حين يعمل المنفِّذ
عبر متغير ثانٍ له مقياسه الزمني. والكبت الآني كان سيستقر عند
مستوى ثابت؛ وأما التأخير بين الاستنساخ والكبت النووي
(الترجمة، والتثنّي، والفسفرة، والاستيراد) فيدع البروتين
يتجاوز قبل أن تُطفأ المورّثة ويهبط دونها قبل أن تُشغَّل من
جديد، والدور يقارب ضعف مجموع التأخير وزمن هدم البروتين. وهدم
PER الأسرع يقصّر الدور: فطفرة مزعزِعة في PER2 البشري تعطي
ساعة نحو 20 ساعة وعائلةً تنام عند السابعة مساءً.
\end{solution}

\begin{solution}{exo:b3:endocrinology:11}
الأس $k(\text{Ca} - 1.2)$: فعند $1.10$، $\mathrm{e}^{-2}$، وPTH $= 1/
(1 + 0.135) = 0.88$؛ وعند $1.15$: $0.73$؛ وعند $1.20$: $0.50$؛ وعند $1.25$: $1/(1 +
2.72) = 0.27$؛ وعند $1.30$: $1/(1 + 7.39) = 0.12$. والميل عند \omterm{thm:b3:endocrinology:loop}{نقطة الضبط}
$-k/4 = -5$ لكل ميلي مول في اللتر: أي \qty{5}{\%} من الأقصى لكل
\qty{0.01}{mmol/L}. والكالسيوم يجب أن يُحفظ في حدود بضع في
المئة ومنفِّذه يعمل على موقٍ ضخم (العظم) بلا تجاوز، فالاستجابة
الحادة آمنة ولازمة. وعروة غلوكوز بهذا الانحدار، تعمل عبر
\omterm{def:b3:endocrinology:glucose}{أنسولين} له زمن تصفيته، كان سيكون لها كسب عروة كبير وتتذبذب
عميقًا في نقص سكر الدم بعد كل وجبة.
\end{solution}

\begin{solution}{exo:b3:endocrinology:12}
التسكّر لا رجعة فيه وبطيء، فتكون في خلية عمرها $\tau$ نسبة
$kG\tau$ مسكَّرة؛ وبالتوسيط على خلايا موزَّعة الأعمار
\numrange{0}{120} يومًا، يكون HbA1c $= kG\times\qty{60}{d}$، متناسبًا
مع متوسط الغلوكوز. والمعايرة: $5.5\times 60k = 0.05$ تعطي $k =
\qty{1.5e-4}{}$ لكل (ميلي مول في اللتر) في اليوم؛ وعند \qty{10}{mmol/L}: \qty{9.1}{\%}.
ومع خلايا تعيش 60 يومًا يكون متوسط العمر 30 يومًا وينتصف
HbA1c للغلوكوز نفسه: فمريض عند \qty{10}{mmol/L} كان سيقرأ
\qty{4.5}{\%}، وهي سوية.
\end{solution}

\begin{solution}{pb:b3:endocrinology:1}
\textbf{1.} \qty{90}{mg/dL} $= \qty{0.9}{g/L} = 0.9/180 = \qty{5.0}{mmol/L}$؛
و$0.9\times 15 = \qty{13.5}{g}$ في حجم التوزع.
\textbf{2.} $75/15 = \qty{5}{g/L}$: أي ارتفاع قدره \qty{500}{mg/dL}،
أي \qty{28}{mmol/L}. والذروة الحقيقية أدنى بكثير لأن الامتصاص
ينتشر على ساعة بينما يزيل التخلص الغلوكوز وهو يصل، ولأن
الكبد يلتقط حصة كبيرة في المرور الأول.
\textbf{3.} $\qty{75}{g}/\qty{60}{min} = \qty{1.25}{g/min}$؛ وعلى
\qty{15}{L}: \qty{8.3}{mg/dL/min}، أي أربعة أمثال $u_{0}$.
\textbf{4.} الفائض المستقر $u/a = 8.3/0.02 = \qty{417}{mg/dL}$، وثابت
الزمن $1/a = \qty{50}{min}$؛ وبعد ساعة $417(1 -
\mathrm{e}^{-1.2}) = \qty{291}{mg/dL}$ فوق الصيام: أي نحو
\qty{380}{mg/dL}، و\qty{21}{mmol/L}.
\textbf{5.} $L = 0.2\times 0.05/(0.02\times 0.1) = 5$؛ والفائض المستقر
$417/6 = \qty{69}{mg/dL}$، أي أصغر بستة أضعاف مما بلا \omterm{def:b3:endocrinology:glucose}{أنسولين}.
\textbf{6.} يحرق الدماغ \qty{120}{g} من الغلوكوز في اليوم، ولا
يخزن منه شيئًا، ولا يقدر على استعمال غيره إلا قليلًا: فغلوكوز
منخفض يعطي تشوشًا وسُباتًا في دقائق. وغلوكوز مرتفع يسكّر
البروتينات ويؤذي، على مدى سنين، أوعية \omterm{def:b3:sensory-systems:retina}{الشبكية} والكلية والعصب
الصغيرة والشرايين الكبيرة.
\textbf{7.} $\lambda^{2} + (a + \gamma)\lambda + (a\gamma + s\beta) = 0$:
والأثر $-0.12$، والمحدد $0.002 + 0.01 = 0.012$، والمميز
$0.0144 - 0.048 = -0.0336$.
\textbf{8.} $\omega = \sqrt{0.0336}/2 = \qty{0.092}{min^{-1}}$؛ والدور
\qty{68}{min}؛ وزمن الخمود $2/0.12 = \qty{16.7}{min}$.
\textbf{9.} $\dot g(0) = 100(-0.06 + c\,\omega) = -2$، فيكون $c\,\omega =
0.04$، و$c = 0.44$.
\textbf{10.} يكون $g = 0$ حين $\tan\omega t = -1/c = -2.29$، أي $\omega t =
\pi - 1.16 = 1.98$، و$t = \qty{21.6}{min}$. وعند أول نهاية صغرى،
قرب \qty{32}{min}: $100\,\mathrm{e}^{-1.94}(\cos 2.97 + 0.44\sin 2.97) =
14.4\times(-0.91) \approx -\qty{13}{mg/dL}$، فالغلوكوز \qty{77}{mg/dL}.
\textbf{11.} عند $t = 0$ يكون $\dot i = 100\beta > 0$: فيرتفع \omterm{def:b3:endocrinology:glucose}{الأنسولين}
ما دام $g > 0$، ويبلغ ذروته حين $\beta g = \gamma i$، وهو ما يقتضي
أن يكون $g$ موجبًا بعد --- فالذروة تأتي قبل أن يبلغ الغلوكوز
الصيام، وبعد أن يبدأ الغلوكوز بالهبوط (وهو يهبط من $t = 0$).
وفي النموذج تقع الذروة قرب \qty{12}{min}.
\textbf{12.} مع دخل منتشر على ساعة تقع الذروة حين يتوازن
الامتصاص والتخلص، أدنى وأمتأخر (30--60 دقيقة)؛ وينتظر الرجوع
انتهاء الامتصاص، ومن ثم ساعتان؛ وتجاوز \omterm{def:b3:endocrinology:glucose}{الأنسولين} أصغر لدخل
تدريجي، فيكون الهبوط خفيفًا.
\textbf{13.} $L = 2.5$. والسليم $g^{*} = (2/0.02)/6 = \qty{16.7}{mg/dL}$؛
والمقاوم $100/3.5 = \qty{28.6}{mg/dL}$: أي ارتفاع قدره نحو
\qty{12}{mg/dL} (\qty{0.7}{mmol/L})، وغلوكوز صيام قرب
\qty{102}{mg/dL}، و\qty{5.7}{mmol/L} --- وهي عتبة اختلال غلوكوز
الصيام.
\textbf{14.} $i^{*} = \beta g^{*}/\gamma$: فالسليم $8.3$، والمقاوم
\qty{14.3}{mU/L}، والنسبة $1.7$: أي مقاومة \omterm{def:b3:endocrinology:glucose}{أنسولين} معوَّضة،
وفرط \omterm{def:b3:endocrinology:glucose}{أنسولين} الدم مع غلوكوز قريب من السويّ.
\textbf{15.} $L = 0.1\times 0.01/0.002 = 0.5$؛ و$g^{*} = 100/1.5 =
\qty{66.7}{mg/dL}$، أي \qty{50}{mg/dL} فوق الحالة السليمة:
و\qty{2.8}{mmol/L}، فغلوكوز الصيام \qty{7.8}{mmol/L} (\qty{140}{mg/dL})،
فوق العتبة السكرية.
\textbf{16.} المميز $(a - \gamma)^{2} - 4s\beta = 0.0064 - 0.004 =
0.0024 > 0$: فالرجوع رتيب. و$\lambda = -0.06 \pm 0.0245$: أي $-0.0355$
و$-0.0845$ في الدقيقة؛ وأبطأ ثابت زمن \qty{28}{min}، في مقابل
\qty{16.7}{min} للخمود السليم: فالعروة المخفقة ترجع أبطأ ولا
تهبط دون المستوى أبدًا.
\textbf{17.} غلوكوز مرتفع مع \omterm{def:b3:endocrinology:glucose}{أنسولين} مرتفع: أي مقاومة بتعويض
جزئي --- النمط~2؛ وقيمة الساعتين \qty{13}{mmol/L} تتجاوز
\qty{11.1}{}، فهو سكري لا مجرد اختلال تحمل.
\textbf{18.} غياب الببتيد C يعني غياب \omterm{def:b3:endocrinology:glucose}{الأنسولين} الداخلي:
فالنمط~1. ومن دون \omterm{def:b3:endocrinology:glucose}{أنسولين} تعجز العضلة والشحم عن التقاط
الغلوكوز، ويُهدَم الشحم إلى أحماض دهنية وكيتونات، ويُهدَم
بروتين العضلة لتوليد الغلوكوز، ويُفقد الغلوكوز في البول
بسعراته: فيجوع المريض في وسط الوفرة.
\textbf{19.} $5\times 9/5.5 = \qty{8.2}{\%}$.
\textbf{20.} $1.5\,\mathrm{e}^{0.46\times 3} = 1.5\times 3.97 =
\qty{6.0}{mU/L}$. فT$_{4}$ ضمن المجال، وTSH فوقه: أي قصور درق
دون سريري، والغدة تخفق \omterm{def:b3:endocrinology:axes}{والنخامى} تعوّض.
\textbf{21.} $1.5\,\mathrm{e}^{0.46\times 8} = 1.5\times 39.6 =
\qty{59}{mU/L}$. وT$_{4}$ منخفض مع TSH لا يتجاوز $0.3$ يعني أن
\omterm{def:b3:endocrinology:axes}{النخامى} لا تستجيب: فالآفة مركزية (نخامية أو وطائية)، لا في
الدرق.
\textbf{22.} عمر نصف واحد: فيهبط المستوى إلى \qty{50}{\%}، وتتسلل
الأعراض على مدى أسابيع. وأما الكورتيزول، بعمر نصف
\qty{80}{min}، فيزول في ساعات من جرعة فائتة، ولا يلقى كربُ
ذلك اليوم دفاعًا: فأسبوع الثيروكسين الفائت مزعج، ويوم
الكورتيزول الفائت قد يقتل.
\textbf{23.} T$_{4}$ مرتفع (\omterm{def:b3:adaptive-immunity:antibody}{فالضد} يسوق الغدة بصرف النظر عن
التغذية الراجعة)؛ وTSH مكبوت بفعل T$_{4}$ المرتفع؛ والعلاقة
اللوغاريتمية الخطية تدفعه إلى ما لا يُكشف --- فعند
$T_{4} = 30$، $1.5\,\mathrm{e}^{-6.9}
= \qty{0.0015}{mU/L}$ --- فيكون TSH غير القابل للكشف أول علامة.
\textbf{24.} مستوى \omterm{def:b3:endocrinology:hormone}{الهرمون} يخبر بميزان أمره وغدته: فTSH مرتفع
مع T$_{4}$ منخفض غدةٌ مخفقة، وTSH مرتفع مع T$_{4}$ مرتفع
\omterm{def:b3:endocrinology:axes}{نخامى} جامحة؛ \omterm{def:b3:endocrinology:glucose}{وأنسولين} مرتفع مع غلوكوز مرتفع جسدٌ مقاوم،
\omterm{def:b3:endocrinology:glucose}{وأنسولين} منخفض مع غلوكوز مرتفع جزيرةٌ مهدَّمة. وT$_{4}$ قدره
\qty{7}{pmol/L} أو \omterm{def:b3:endocrinology:glucose}{أنسولين} قدره \qty{30}{mU/L} لا يقول، مقروءًا
وحده، أين الخلل؛ ومقروءًا مع آمره، يقول.
\textbf{25.} \omterm{thm:b3:endocrinology:loop}{كسب العروة} $L = 5$؛ والدور \qty{68}{min} وزمن الخمود
\qty{17}{min}؛ وغلوكوز الصيام في العروة المخفقة \qty{7.8}{mmol/L}.
\end{solution}
